% James's text, copied from his draft with the wording unchanged
% only latex escapes, citation commands, and table references were added
% replace this file with the final version from the doc before submitting

How do different LLMs perform on cross-domain Text-to-SQL? This research question was approached and answered with a zero-shot benchmark of 11 different configurations of LLMs on the full Spider 1.0 development dataset, with the same prompt, metrics, and statistics being used for each configuration.

\paragraph{Models.}
We chose these 11 configurations to benchmark by manipulating three different factors: family, size and inference mode. The first family we chose was Qwen 2.5-Coder, a decoder-only LLM (no explicit reasoning), and within that family we varied the size at the following increments: 0.5B, 1.5B, 3B, and 7B. The second family we chose was Qwen3, a reasoning LLM, and within that family we varied the size three times at 1.7B, 4B, and 8B. Qwen3 is a hybrid model, meaning it supports both thinking and non-thinking, and so we were able to alter its inference mode, using the same checkpoints, and without altering the weights \citep{yang2025qwen3}. Thus we obtained three more Qwen3 configurations, one at each aforementioned size, to compare against the adjacent thinking configurations. Additionally we included a 4-bit quantized configuration of Qwen2.5-Coder-7B with AWQ to assess the accuracy and speed cost of weight quantization \citep{lin2024awq}. Throughout this report we will be referring to these three factors that discern between and describe the configurations: family, or the training recipe (e.g.\ Qwen2.5-Coder or Qwen3), size, or the model’s parameter count (measured in billions), and specifically for Qwen3, the inference mode, or whether thinking is enabled or disabled for that configuration.

All eight of the checkpoints used have a dense decoder-only transformer architecture. Although Qwen3 is a reasoning LLM,  its dense models do share a very similar architecture to Qwen2.5 models \citep{yang2025qwen3}. Therefore the Qwen3 models major differences from the Qwen2.5-Coder models come from how they were trained, rather than their network architectures. Specifically Coder is Qwen2.5 trained on about 5.5 trillion tokens of mostly code, making it a code specialized model with no explicit reasoning, while Qwen3 was trained with reinforcement learning for reasoning, merged into a single model that supports both modes and smaller sizes are derived from larger Qwen3 models \citep{hui2024qwen25coder,yang2025qwen3}. Since none of the models use a Mixture of Experts (MoE) design, the total and active parameter counts were identical for each checkpoint. Holding the broad architecture relatively constant allowed us to focus comparisons on model family, model size and inference mode. However, these factors were not equally isolated. The most isolated comparison was between the thinking and non-thinking inference modes, as both modes run the exact same Qwen3 checkpoint at each size, making the thinking setting in the chat template the only difference. Comparisons across parameter sizes were also largely isolated, as models in the same family share the same training data and pipeline, with scaling to larger models primarily using more and wider layers. Comparing across families was the least isolated, as these families differ in their training and their training data, and their sizes do not match exactly. Although this comparison targets the effect of training, the combination of these differences means that a performance difference between the families cannot be attributed to any single factor. Instead, from this comparison we can reveal how each training recipe performs as a whole, which still answers if a code specialized model or a more general model trained to reason performs better on a Text-to-SQL task. Furthermore, comparing Qwen2.5-Coder to Qwen3 with thinking disabled, and then against Qwen3 with thinking enabled can help separate the differences in performance due to reasoning itself, and the rest of the differences between the families.

All of the above configurations were run unquantized in bfloat16, our one extension from this was the quantized configuration that used Qwen’s official AWQ release of Qwen2.5-Coder-7B-instruct \citep{qwen2024awqcard}. AWQ (Activation aware Weight Quantization) distinguishes roughly the top 1\% of weight channels that matter the most within a model by looking at their size of activations passing through them and protects them from large quantization errors by scaling them up prior to quantization \citep{lin2024awq}. After scaling was applied the model weights were then quantized to 4-bit, while computation runs in float16. This extension was done to see if reducing the necessary amount of memory would improve speed without introducing too much loss of accuracy. Therefore, including both the bfloat16 and 4-bit AWQ versions of Qwen2.5-Coder-7B allowed for the assessment of this tradeoff between efficiency and accuracy.

All inference was performed locally. The weights were downloaded from the Hugging Face Hub at pinned revisions in October 2026 and were served with a vLLM, no API was used. Table~\ref{tab:models} shows the exact Hugging Face identifier, revision, parameter count, and precision for each checkpoint, with parameter counts taken from each model’s card.

\begin{table}[htbp]
\centering
\caption{Model Configurations evaluated in Task 1}
\label{tab:models}
\scriptsize
\setlength{\tabcolsep}{3pt}
\resizebox{\linewidth}{!}{\begin{tabular}{@{}llllllll l@{}}
\toprule
Model (Hugging Face ID) & Family & Params & Non-emb. & Modes & Decoding & Max new tok. & Precision & Rev. \\
\midrule
\texttt{Qwen/Qwen2.5-Coder-0.5B-Instruct} & Coder & 0.49B & 0.36B & standard & greedy & 512 & bf16 & \texttt{ea3f247} \\
\texttt{Qwen/Qwen2.5-Coder-1.5B-Instruct} & Coder & 1.54B & 1.31B & standard & greedy & 512 & bf16 & \texttt{2e1fd39} \\
\texttt{Qwen/Qwen2.5-Coder-3B-Instruct} & Coder & 3.09B & 2.77B & standard & greedy & 512 & bf16 & \texttt{488639f} \\
\texttt{Qwen/Qwen2.5-Coder-7B-Instruct} & Coder & 7.61B & 6.53B & standard & greedy & 512 & bf16 & \texttt{c03e6d3} \\
\texttt{Qwen/Qwen2.5-Coder-7B-Instruct-AWQ} & Coder & 7.61B & 6.53B & standard & greedy & 512 & 4 bit weights, fp16 & \texttt{8e8ed24} \\
\midrule
\texttt{Qwen/Qwen3-1.7B} & Qwen3 & 1.7B & 1.4B & off / on & greedy / sampled & 512 / 8,192 & bf16 & \texttt{70d244c} \\
\texttt{Qwen/Qwen3-4B} & Qwen3 & 4.0B & 3.6B & off / on & greedy / sampled & 512 / 8,192 & bf16 & \texttt{1cfa9a7} \\
\texttt{Qwen/Qwen3-8B} & Qwen3 & 8.2B & 6.95B & off / on & greedy / sampled & 512 / 8,192 & bf16 & \texttt{b968826} \\
\bottomrule
\end{tabular}
}

\vspace{2pt}
\begin{minipage}{\linewidth}\scriptsize
\textit{Note.} Parameter counts from the model cards (Qwen3 cards round to one decimal); sampled = temperature 0.6, top-p 0.95, top-k 20; thinking seeds 42, 66, 73 (plus 137, 255 for 4B); Rev. = first 7 characters of the Hugging Face commit hash.
\end{minipage}
\end{table}

\paragraph{Evaluation data.}
All evaluation data was performed on the development set of the Spider 1.0 dataset containing 1,034 natural language questions over 20 databases \citep{yu2018spider}. Since this dataset is a cross-domain dataset the development databases do not appear in the training split. Therefore, the models must be able to generalize to database schemas not represented by the Spider training examples rather than answering new questions from databases seen in the training split. Our benchmark is zero-shot, so the training split was not used for either few-shot examples and it also wasn’t used for fine tuning or adaptation. Furthermore, we did not use a subsample of the development split, instead every model configuration was evaluated on the whole set of 1,034 development questions. This decision removes the potential for sampling to bias specific models or difficulty levels. The project seed of 42 was fixed prior to running any experiments but had no effect on which questions were being evaluated since the whole split was used. The difficulty of the questions was determined by using the spider’s official hardness function script and assigned each question a difficulty level of easy, medium, hard, or extra hard depending on the properties of its gold SQL query. The majority of the gold queries in the development set appear twice and are paired with two differently worded questions. The total of 1,034 questions actually corresponds to 551 gold queries. However, these questions that correspond to the same gold query and not fully independent observations and so our statistical analysis must account for this pairing.

\paragraph{Prompting and input.}
For prompting and input we used the same prompt for every configuration and it was fixed before any results were seen. Every model received the same underlying zero-shot prompt as to limit any additional source of variation. The prompt itself was made up of two messages, a system message and a user message. First the system message would say: “You are an expert in SQLite. Given a database schema and a question, write one SQLite query that answers the question”.The user message would give the database schema under the header “\#\#\# Database schema”, followed by the question under the header “\#\#\# Question”, and ended with the single instruction: “Return only the final SQL query inside a sql code block”. Each database schema was given with the databases’s CREATE TABLE statements. The statements can then be read directly from each SQLite file, exactly as it was stored, including column table names, column names, column types, and primary and foreign keys. \citet{rajkumar2022evaluating} found that this method of using CREATE TABLE statements was effective in representing relational schemas to LLMs. However, their strongest version provided three example rows for each table. We intentionally did not utilize this exact same format in order to keep the prompts shorter, and give all of the models the same purely schema information, which prevents database contents from becoming an additional source of information. This denial of example values and content does also limit useful information and may reduce overall performance, and therefore is a limitation of the evaluation. The final instruction of the prompt requesting only the final SQL query inside a sql block aided in our extraction of the final SQL from each model’s output. For Qwen2.5-Coder and Qwen non-thinking models, the output should be the answer. For Qwen3 thinking models the procedure is different because the output can contain a reasoning trace followed by the final SQL. In this case Qwen3 the output was split at the closing think token, and only the text after that token is searched for SQL, thus all queries during reasoning are ignored. If the think token never occurs we conclude that the entirety of the output token allowance was consumed prior to finishing and it is recorded as having no prediction. After separating the answer from the reasoning trace, the same extraction procedure was applied to all configurations. First the primary extraction searched for the last fenced code block to be taken as prediction. If there was no code block present then a regular expression was used to extract the text from the first SELECT or WITH as in spider the model’s predicted SQL query typically begins with one of these keywords. If neither of these methods produced a SQL query, the output was then recorded as having no prediction and scored incorrect. Across all the 19,646 outputs collected from single runs, 19,620 predictions were extracted directly from SQL code blocks, another 10 required the regular expression, and 16 had no prediction. In all 16 cases the reasoning trace reached the maximum output token allowance. The full prompt template and an example prompt are given in Appendix~\ref{app:prompt}.

\paragraph{Inference settings.}
As shown in Table~\ref{tab:inference} the inference settings were divided into non-thinking and thinking configurations. The non-thinking models used greedy decoding while the thinking models used sampling. The non-thinking models used greedy decoding because the model always selects the highest probability next token. That makes this method of decoding deterministic, so with the same model and prompt, the resulting SQL prediction should always be the same. This is why using the standard seed of 42 doesn’t change anything. This allows us to have a deterministic benchmark that can be used for clean size comparisons within families, and more consistent comparisons between families by removing sampling variability from the comparisons. The other settings include a temperature of 0, a top-p of 1, and top-k off with up to 512 new tokens. Since temperature, top-p, and top-k are all controls for influencing sampling they are not required, and 512 tokes gives ample space for a SQL query. The thinking models use sampling instead as it allows for the model to select the next reasoning step from a pool of high probability options, thus the same question with the same Qwen3 weights, but a different seed could produce a different reasoning trace and potentially a different SQL. The specific settings of temperature = 0.6, top-p = 0.95, and top-k = 20 are directly from the official Qwen3 model card \citep{qwen2025qwen3card,yang2025qwen3}. We followed the model card for all settings except for Qwen3 greedy without thinking, and the max token cap for Qwen3 thinking. We changed the Qwen3 greedy without thinking from temperature = 0.7, top-p = 0.8, and top-k = 20, to the greedy decoding settings above to make those runs deterministic and directly comparable to the performance of the Coder runs. We also reduced the max token cap from 32,768 to 8,192 out of necessity. The vLLM’s context window had to hold the prompt and the output, and while it was set to 16,384, the largest model, or the 8B parameter model reduced it to a 10,240 token window, to fit in the available GPU memory. Since the longest prompt was roughly 1,000 tokens, we had to cap it close to 8,000 tokens. However, this cap was only reached in 16 of 11,374 thinking outputs, and between 0 to 3 per run, so this did not negatively impact the performance. For thinking models seeds 42, 66, and 73 were used, and additionally for 4B 137 and 255 were used. For Qwen3 models specifically, reasoning was switched on or off using the model’s chat template, as shown in Table~\ref{tab:inference}. We ran these models through Google Colab where we used a single NVIDIA L4 GPU with 24 GB VRAM using vLLM 0.30.0 \citep{kwon2023vllm}, and the unquantized model configurations loaded in bfloat16 precision. The questions were processed in batches, each of 100 prompts, and vLLM generated every prompt in a batch together. Seconds per question was then calculated by dividing the total inference time of each batch by 100. Software versions and full revision hashes for every run are listed in Appendix~\ref{app:software}.

\begin{table}[htbp]
\centering
\caption{Inference settings by mode}
\label{tab:inference}
\small
\begin{tabular}{@{}lll@{}}
\toprule
Setting & Coder and Qwen3 no thinking & Qwen3 thinking \\
\midrule
Decoding & greedy & sampling \\
Temperature / top-p / top-k & 0 / 1.0 / off & 0.6 / 0.95 / 20 \\
Max new tokens & 512 & 8,192 \\
Reasoning & off & on (\texttt{enable\_thinking}) \\
Seeds & 42 & 42, 66, 73 (plus 137, 255 for 4B) \\
Context window & 16,384 (10,240 for 8B) & 16,384 (10,240 for 8B) \\
\bottomrule
\end{tabular}


\vspace{2pt}
\begin{minipage}{0.92\linewidth}\footnotesize
\textit{Note.} Each Qwen3 checkpoint appears in both columns, so the only differences between its two modes are the ones listed here. Because greedy decoding always returns the same output, seed 42 has no effect in the left column. Thinking runs sample, so their results vary by seed and were repeated. The context window is the total prompt plus output length vLLM allowed, and it was reduced for the 8B so the model fit in memory.
\end{minipage}
\end{table}

\paragraph{Evaluation.}
\todo{Write: metrics and why (EM and component F1 are what the Spider paper reports; EX with the model's own values is primary; EX with gold values is the paper's own EX definition and an upper bound); how they were computed (official test-suite-sql-eval, single Spider database, defaults: row order ignored unless ORDER BY, DISTINCT removed, value-free EM; per-question scores match the official totals on every run); invalid and unparseable outputs (denominator 1,034; no prediction scores 0 on both; crash scores 0 on EX; 2018 parser failure scores 0 on EM but EX still runs; 30 s limit, 1 timeout in 19,646; gold-value cap of 1,000 combinations or about 120 s, 28 of 19,646 capped); statistics (Wilson 95\% intervals; 13 planned comparisons with exact McNemar and Holm within each metric; bootstrap over the 551 gold queries, 2,000 resamples, seed 42; seed SD; pessimistic EX on the 73 trivial gold answers). Component F1 is in Appendix~\ref{app:components} (Table~\ref{tab:components}). Cite \citet{zhong2020testsuite}, \citet{mcnemar1947}, \citet{dietterich1998}, \citet{wilson1927}, \citet{holm1979}, \citet{miller2024errorbars}, \citet{kim2025flex}.}
